\chapter{Черные дыры как естественные лаборатории для исследования сильных гравитационных полей}\label{ch:1}
Спустя более чем 100 лет своего существования Общая Теория Относительности, созданная Эйнштейном, является одной из наиболее успешных физических теорий, предсказания которой были подтверждены большим количеством экспериментов и наблюдений. Несмотря на это,  далеко не все аспекты теории хорошо изучены: сверхмассивные компактные объекты, полностью гравитационно сколлапсированные объекты или черные дыры в настоящий момент остаются малоизученными объектами. Лишь недавно удалось непосредственно пронаблюдать тень сверхмассивной черной дыры в центре галактики — объект, предсказанный более ста лет назад.

Предположение о том, что черные дыры могут быть обнаружены при помощи наземных наблюдений, получило большое развитие после открытия нейтронных звезд (пульсаров) в конце 1960-х годов. Первый серьезный кандидат на роль черной дыры, Лебедь X-1, обнаружен в 1972 году Чарльзом Т. Болтоном, Луизой Вебстер и Полом Мердином \cite{boltonDimensionsBinarySystem1972, websterCygnusX1aSpectroscopic1972} с помощью рентгеновского излучения двойной звездной системы, обычно состоящей из тяжелой звезды (чаще черной дыры), которая постепенно наращивает массу за счет более легкой звезды-компаньона. Примерно через год после этого было принято во внимание, что гипотеза о черной дыре дает простейшее объяснение наблюдаемому рентгеновскому спектру, хотя некоторое время оставались сомнения относительно того, содержит ли Лебедь X-1 черную дыру. В то же время с помощью рентгеновских спектров двойных систем были обнаружены десятки черных дыр \cite{tetarenkoWATCHDOGCOMPREHENSIVEALLSKY2016}.

В 1994 году Чарльз Таунс и его коллеги \cite{genzelNucleusOurGalaxy1994, danchiCharacteristicsDustShells1994} наблюдали сильно вращающийся ионизированный неоновый газ вокруг центра нашей галактики в созвездии Стрельца А*. Они приписали свое наблюдение сверхмассивной черной дыре массой около нескольких миллионов масс Солнца. В настоящее время мы имеем прямые наблюдения аккреционного диска этой черной дыры при помощи Event Horizon Telescope (EHT). 

%Сложно предсказать дальнейшее развитие технологий, но можно надеяться на то, что в ближайшие десятилетия качества получаемых изображений теней черных дыр значительно улучшится и появится возможность проверки нашей теории гравитации на новых наблюдательных данных.

\section{Черные дыры как точные решения уравнений Эйнштейна}\label{sec:1.1}
\subsection{Исторические предпосылки и первые решения}\label{sub:1.1.1}
Идея объектов, обладающих столь сильной гравитацией, что даже свет не может их покинуть, впервые возникла еще в рамках ньютоновской механики. В 1783 году английский естествоиспытатель Джон Мичелл, используя корпускулярную теорию света Ньютона и закон всемирного тяготения, рассчитал, что для массивного тела с радиусом, превышающим критическое значение (R > 2GM/c²), хотя точная формулировка появилась позже), вторая космическая скорость будет больше скорости света. Независимо от него в 1796 году французский математик Пьер-Симон Лаплас пришел к аналогичным выводам в своей работе "Exposition du système du Monde". Эти гипотетические объекты получили название "темных звезд" (dark stars) и рассматривались как теоретически возможные, хотя их физическая природа оставалась неясной в отсутствие релятивистской теории гравитации.

Несмотря на кажущуюся схожесть с современной концепцией черных дыр, "темные звезды" ньютоновской физики принципиально отличались от релятивистских объектов: в ньютоновском подходе свет рассматривался как частицы, подчиняющиеся механическим законам, а не как электромагнитные волны с постоянной скоростью в вакууме. Кроме того, отсутствовало понимание таких ключевых релятивистских эффектов как искривление пространства-времени, гравитационное красное смещение и горизонт событий. Однако, работы Мичелла и Лапласа представляют исторический интерес как первые попытки осмыслить предельные проявления гравитации, предвосхитившие более поздние открытия в общей теории относительности.

Формулировка общей теории относительности (ОТО) Альбертом Эйнштейном в 1915 году создала принципиально новую парадигму описания гравитации через искривление пространства-времени, что позволило корректно описать отклонение траектории движения света вблизи Солнца от прямолинейной и объяснить аномальную прецессию перигелия Меркурия. Основой данного подхода является описание пространства при помощи метрического тензора (или метрики) \(g_{\mu \nu}\), который определяет интервал \(ds\) между двумя бесконечно близкими событиями:
\begin{equation}
    \label{eq:general_interval}
    ds^2=g_{\mu \nu}dx^{\mu}dx^{\nu},
\end{equation}
где \(dx^{\mu}\) — 4-вектор, соединяющий два бесконечно близких события в некоторой системе координат.
Метрика полностью определяет траектории движения тел, являющиеся геодезическими в искревленном пространстве времени. Сам же метрический тензор зависит от распределения материи в пространстве и может быть найден путем решения уравнения Эйнштейна:
\begin{equation}
    \label{eq:Ein_Eq}
    R_{\mu \nu}-g_{\mu \nu}R=8\pi G T_{\mu \nu},
\end{equation}
где \(R_{\mu \nu}\) — тензор Ричи, получаемый из метрики, \(R=R_{\mu \nu}g^{\mu \nu}\) — скаляр кривизны, \(T_{\mu \nu}\) — тензор энергии-импульса материи, \(G\) — гравитационная постоянная Ньютона. В дальнейшем для упрощения выкладок во всех формулах  будет принята система единиц, в которой \(G=c=k=1\), где \(с\) — скорость света в вакууме, \(k\) — постоянная Больцмана.

В 1916 году Карлом Шварцшильдом найдено первое точное решение уравнений Эйнштейна для сферически-симметричного вакуумного случая \cite{schwarzschild1916gravitationsfeld}:
\begin{equation}
    \label{eq:Schw_Metric}
    ds^2=(1-\frac{2M}{r})dt^2 - (1-\frac{2M}{r})^{-1}dr^2 - r^2d\theta^2 - r^2\sin{\theta}^2d\phi^2
\end{equation}
где \(M\) — масса искривляющей пространство материи.

Данная метрика содержит две ключевые особенности: горизонт событий при \(r = 2M\) и сингулярность при \(r = 0\), что стало математическим прообразом современного понятия черной дыры. Это решение заложило основы дальнейшего изучения компактных релятивистских объектов.

Почти одновременно с работой Шварцшильда началось исследование более сложных уравнений, учитывающих электромагнитное поле. В 1916-1918 годах Ханс Райсснер \cite{reissner1916eigengravitation} и Гуннар Нордстрём \cite{nordstrom1918energy} независимо получили решение для сферически-симметричного заряженного тела:
\begin{equation}
    \label{eq:Reis-Nord_Metric}
    ds^2=(1-\frac{2M}{r} +\frac{Q^2}{r^2})dt^2 - (1-\frac{2M}{r}+\frac{Q^2}{r^2})^{-1}dr^2 - r^2d\theta^2 - r^2\sin{\theta}^2d\phi^2
\end{equation}
Полученная метрика содержит дополнительный член, описывающий вклад электрического заряда \(Q\). Это решение продемонстрировало существование внутреннего горизонта Коши и сложную причинную структуру, что впоследствии привело к формулировке проблемы инфляции массы.

\nomenclature{\(M_{\bigodot}\)}{Масса Солнца}

В 1963 году математик Рой Керр в работе \cite{kerr1963gravitational} представил точное решение для стационарной аксиально-симметричной черной дыры с угловым моментом \(J\). Полученная им метрика в координатах Бойера-Линдквиста имеет вид
\begin{equation}
	\label{eq:Kerr-New_Metric}
	ds^2=-\frac{\Delta}{\Sigma}[dt-a\sin^2\theta d\phi]^2+\frac{\sin^2\theta}{\Sigma}[(r^2+a^2)d\phi -adt]^2  +\frac{\Sigma}{\Delta}dr^2+\Sigma d\theta^2\!,
\end{equation}
\begin{subequations}
    \label{eq:Kerr_DS}
    \begin{gather}
        \Delta_K = r^2-2Mr+a^2 \label{eq:Kerr_D} \\
        \Sigma= r^2+a^2\cos^2\theta \label{eq:Kerr-New_S}
    \end{gather}
\end{subequations}
где \(a=\frac{J}{M}\) — угловой момент единицы массы.
Это решение, описывающее так называемую "эргосферу" и эффект увлечения инерциальных систем отсчета (frame-dragging), стало основой современной астрофизики черных дыр.

В 1965 году в работах Эца Ньюмена \cite{newmanNoteKerrSpinningParticle1965, newmanMetricRotatingCharged1965} было получено наиболее общее решение, описывающее черную дыру с массой \(M\), зарядом \(Q\) и угловым моментом \(J\) и объединяющее все предыдущие решения. Метрика в этом случае выглядит так же как метрика Керра \cref{eq:Kerr-New_Metric}, отличаясь от нее функцией \(\Delta\):
\begin{equation}
    \label{eq:Kerr-New_D}
    \Delta_{KN} = r^2-2Mr+a^2+Q^2
\end{equation}

\subsection{Формирование современной концепции черных дыр}\label{sub:1.1.2}
Современное понимание черных дыр как физических объектов сформировалось в середине XX века. В 1939 году Роберт Оппенгеймер и Хартланд Снайдер \cite{oppenheimerContinuedGravitationalContraction1939} впервые рассмотрели процесс гравитационного коллапса массивной звезды в рамках ОТО, показав, что при определенных условиях коллапс становится необратимым. Их работа, выполненная через несколько лет после открытия нейтронов, впервые связала теоретические решения типа Шварцшильда с реальными астрофизическими процессами. Однако сам термин "черная дыра" был введен Джоном Арчибальдом Уилером в 1967 году, хотя аналогичные понятия ("застывшая звезда", "коллапсар") использовались и ранее советскими физиками.

Важнейший вклад в формирование строгой теории черных дыр внес Роджер Пенроуз в 1965 году \cite{penroseGravitationalCollapseSpaceTime1965a}, доказав первую теорему о сингулярностях, которая устанавливала неизбежность образования сингулярности при гравитационном коллапсе в рамках ОТО. Его работа ввела принципиально новые математические методы (теория глобальной причинности, диаграммы Пенроуза) для анализа решений Эйнштейна. В 1971 году Стивен Хокинг \cite{SingularitiesGravitationalCollapse1970} обобщил этот результат, показав, что сингулярности должны возникать и в космологических моделях. Пенроуз также сформулировал гипотезу "космической цензуры" (1969) \cite{penroseGravitationalCollapseRole1969}, утверждающую, что все сингулярности в реалистичных условиях должны быть скрыты за горизонтом событий.

Революционный поворот в понимании черных дыр произошел в 1970-х годах с открытием их термодинамических свойств. Яков Бекенштейн в 1972 году \cite{bekensteinBlackHolesEntropy1973} предположил, что черные дыры обладают энтропией, пропорциональной площади горизонта \(S = \frac{A}{4}\). Это утверждение первоначально вызвало сопротивление, пока Стивен Хокинг в 1974-1975 годах \cite{hawkingParticleCreationBlack1975} не показал, что черные дыры действительно должны испускать квантовое излучение с температурой \(T = \frac{\hbar}{8\pi M}\) (излучение Хокинга). Эти работы заложили основы квантовой теории черных дыр и связали гравитацию, термодинамику и квантовую теорию поля.

Современная концепция черных дыр как фундаментальных астрофизических объектов окончательно сформировалась к концу XX века благодаря развитию теории аккреции и появлению наблюдательных данных. 
%Открытие квазаров (1963), рентгеновских двойных систем (1970-е) и, наконец, прямое обнаружение гравитационных волн от слияния черных дыр (LIGO, 2015) \cite{abbottObservationGravitationalWaves2016a} и изображение тени черной дыры (Event Horizon Telescope, 2019) \cite{EHTSgrA2022, EHTM872019} превратили математическую абстракцию начала XX века в общепризнанный астрофизический феномен. Сегодня черные дыры рассматриваются как ключевые объекты для проверки ОТО в сильных полях и поиска квантовой гравитации.

%\section{Обнаружение черных дыр}\label{sec:1.2}
%ОТО уже не раз показала свою силу и универсальность. Множество предсказаний, сделанных с использованием ОТО были проверены и подтверждены: космологические модели расширяющейся вселенной, гравитационные волны и, конечно же, черный дыры. Однако до сих пор нам не удалость построить теории квантовой гравитации, которая объединяла бы в себе достижения ОТО и квантовой механики. 
%
%Однако, с развитием технологий нам открывается все больше возможностей по наблюдению таких физических процессов, для которых было бы важно учесть и квантовые и гравитационные эффекты. В частности, вблизи и непосредственно внутри черных дыр должны начать проявляться квантовые эффекты, например излучение Хокинга. В таких экстремальных условиях должно быть проще обнаружить отклонения от предсказаний ОТО, которые бы показали в каком направление необходимо искать новые теории и помогло бы сузить число потенциальных теорий.
%
%Но прежде, чем искать новую теорию, сначала нужно показать, что наблюдения не сходятся с предсказаниями ОТО. Поэтому так важно для начала иметь предсказания для ОТО.
%
%\textbf{Гравитационное линзирование}

%Мы можем сделать \textbf{жирный текст} и \textit{курсив}.

\section{Экспериментальное обнаружение черных дыр}\label{sec:1.2}

Экспериментальное обнаружение черных дыр началось в 1960-х годах с открытием мощных космических рентгеновских источников. Первым кандидатом в черные дыры стал источник Лебедь X-1, обнаруженный в 1964 году при помощи ракетных экспериментов. В 1971 году точные наблюдения с помощью спутника Uhuru позволили отождествить этот источник с голубым сверхгигантом HDE 226868, а анализ орбитального движения показал наличие невидимого компактного объекта массой около 15 \(M_{\bigodot}\) - слишком массивного для нейтронной звезды. Параллельно развивалась теория аккреции (Шакура, Сюняев, 1973), объяснившая высокую светимость таких систем падением вещества на компактный объект. К концу 1970-х были обнаружены десятки рентгеновских двойных систем с кандидатами в черные дыры звездных масс.

В 1970-80-х годах накопились убедительные доказательства существования сверхмассивных черных дыр в центрах галактик: наблюдения за движением газа и звезд в ядре активной галактики M87 (Форд, 1981) и нашей Галактики (Генцель, 1996) показали концентрации массы в миллионы солнечных масс в объемах, исключающих альтернативные объяснения. Особенно убедительными стали данные инфракрасного интерферометра на телескопе VLT, позволившие в 2002 году проследить орбиту звезды S2 вокруг Sgr A* и доказать, что центральный объект массой \(\num{4e6} M_{\bigodot}\)  действительно является черной дырой. В этот же период рентгеновские спутники Chandra и XMM-Newton обнаружили характерные спектральные линии железа, искаженные релятивистскими эффектами вблизи горизонта событий.

Фундаментальный прорыв в экспериментальном изучении черных дыр произошел 14 сентября 2015 года, когда детекторы LIGO зарегистрировали гравитационные волны от слияния двух черных дыр (событие GW150914) с массами 36 и 29 солнечных масс. Этот результат не только подтвердил существование черных дыр как физических объектов, но и продемонстрировал их способность сливаться, образуя более массивные компактные объекты. Последующие наблюдения коллабораций LIGO-Virgo-KAGRA выявили целую популяцию двойных черных дыр, включая экзотические системы с необычными массами и параметрами орбит. Особый интерес представляет событие GW190521 (2020 год) с образованием черной дыры промежуточной массы (~142 \(M_{\bigodot}\)), что заполнило важный пробел в понимании теории фомирования черных дыр.

В апреле 2019 года коллаборация Event Horizon Telescope (EHT) представила первое в истории прямое изображение "тени" сверхмассивной черной дыры M87* - темной области, окруженной асимметричным кольцом излучения, что идеально соответствовало предсказаниям ОТО для керровской черной дыры. Эти наблюдения, проведенные в апреле 2017 года с использованием глобальной сети радиотелескопов, подтвердили существование горизонта событий с беспрецедентной детализацией. В 2022 году EHT опубликовал изображение черной дыры в центре Млечного Пути (Sgr A), демонстрируя удивительное сходство с M87 несмотря на разницу в масштабах. Современные проекты, такие как расширение EHT на космические антенны и разработка новых алгоритмов обработки, обещают еще более детальные изображения и возможность изучать динамику аккреционных процессов в реальном времени.

\section{Верификация ОТО через анализ изображений теней черных дыр}\label{sec:1.3}
\subsection{Современные наблюдения теней черных дыр}\label{sub:1.3.1}
Наблюдение теней сверхмассивных черных дыр представляет собой уникальный инструмент для проверки предсказаний общей теории относительности в режиме сильных гравитационных полей. Теоретический анализ показывает, что форма и размер тени черной дыры (области, где световые лучи захватываются горизонтом событий) чувствительны к параметрам метрики пространства-времени и могут быть точно рассчитаны в рамках ОТО. Первое изображение тени черной дыры M87, полученное коллаборацией Event Horizon Telescope (EHT) в 2019 году, продемонстрировало асимметричное кольцо излучения диаметром \(42±3\) микросекунд дуги, что полностью соответствует предсказаниям для керровской метрики с параметром вращения \(a = 0.90±0.05\). Последующие наблюдения в поляризованном свете (2021) дополнительно подтвердили релятивистскую модель синхротронного излучения в аккреционном диске и структуру магнитных полей, предсказанную GRMHD-симуляциями.

Сравнение наблюдательных данных EHT с теоретическими предсказаниями позволяет установить строгие ограничения на возможные отклонения от ОТО. Анализ формы тени Sgr A* (2022) показал, что наблюдаемый профиль совпадает с керровским решением с точностью лучше \(10\%\), что исключает значительные модификации метрики в рамках таких альтернативных теорий как f(R)-гравитация или теория Хорндески. Современные разработки, включая расширение EHT на более короткие волны и добавление космических антенн, позволят достичь беспрецедентной точности \(~1\%\) в тестировании предсказаний ОТО и поиске возможных квантовых поправок на масштабах, сравнимых с гравитационным радиусом.

\subsection{Моделирование силуэта}\label{sub:1.3.2}

\begin{figure}[ht]
    \centerfloat{
        \hfill
        \subcaptionbox[List-of-Figures entry]{Черная дыра Шварцшильда\label{fig:Shad0_a}}{
            \includegraphics[width=0.45\linewidth]{Shad0_a}}
        \hfill
        \subcaptionbox{Черная дыра Керра\label{fig:Shad0_b}}{
            \includegraphics[width=0.45\linewidth]{Shad0_b}}
        \hfill
    }
%    \legend{}
	\caption[Фотонные траектории и тени черных дыр]{Фотонные траектории и тени черных дыр в разных метриках}\label{fig:Shad0}
\end{figure}

Пурпурная область на небесной сфере – классическая тень черной дыры (сечение захвата фотона черной
дырой). Синий диск – изображение горизонта событий в пространстве Эвклида (без гравитации Многоцветные
кривые — численно вычисленные (геодезические) траектории фотонов.

Удаленный наблюдатель расположен в экваториальной
плоскости черной дыры. Параметрическое уравнение для тени \((\lambda, q) = (\lambda(r), q(r))\)

\begin{subequations}
    \label{eq:1.6}
    \begin{gather}
        \lambda = \frac{(3-r)r^2-a^2(r+1)}{a(r-1)} \label{eq:1.6a} \\
        Q = q^2 = \frac{r^3[4a^2-r(r-3)^2]}{a^2(r-1)^2} \label{eq:1.6b}
    \end{gather}
\end{subequations}
где \((\lambda\) — фотонный горизонтальный и \(q\) – вертикальный прицельный параметр, \(Q\) — постоянная Картера, стрелкой отмечена \(a = 0.65\) — ось вращения черной дыры, пунктирной окружностью — горизонт событий \(r_h = 1+\sqrt{1-a^2}\) \cite{bardeen1973timelike,chandrasekhar1998mathematical,dokuchaevSeeInvisibleImage2019}.

\begin{figure}[ht]
    \centerfloat{
        \hfill
        \subcaptionbox[List-of-Figures entry]{\(a=1\)\label{Shad0_2d_a}}{
            \includegraphics[width=0.4\linewidth]{Shad0_2d_a}}
        \hfill
        \subcaptionbox{\(a=1\)\label{fig:Shad0_2d_b}}{
            \includegraphics[width=0.45\linewidth]{Shad0_2d_b}}
        \hfill
        }
    \caption[Контур тени черной дыры на ярком фоне]{Контур (граница) тени черной дыры на ярком фоне}\label{fig:Shad0_2d}
\end{figure}


Уравнения движения пробных частиц


Визуальные изображения сверхмассивных черных дыр M87* и SgrA* недавно были успешно зарегистрированы международной коллаборацией Event Horizon Telescope [1, 2]. Эта регистрация открывает уникальную возможность для проверки (или опровержения) модифицированных теорий гравитации в пределе сильного поля, где гравитационное поле доминирует над астрофизическими факторами. Такая проверка крайне важна для физической интерпретации астрофизических и космологических наблюдений Вселенной, а также для понимания природы загадочных темной материи и темной энергии.

Точности Event Horizon Telescope достаточно для демонстрации качественного соответствия полученных изображений темных пятен обеих черных дыр предсказаниям Общей теории относительности [3–8]. В частности, форма темных пятен на изображениях Event Horizon Telescope соответствует быстрому вращению сверхмассивных черных дыр M87* и SgrA* [9, 10]. Однако текущая точность Event Horizon Telescope недостаточна для проверки модифицированных теорий гравитации.

Наблюдаемые темные пятна всегда проецируются внутри ожидаемого положения классической тени черной дыры на небесной сфере. С физической точки зрения классическая тень черной дыры представляет собой сечение захвата фотонов в гравитационном поле черной дыры [11–13].

Физически обоснованная модель генерации мощного энерговыделения у аккрецирующих черных дыр основана на механизме Блэнфорда–Знаека [14], в котором через горизонт событий черной дыры протекает электрический ток. Этот ток создает электромагнитный поток с вектором Пойнтинга вдоль оси вращения черной дыры, а также обеспечивает сильный нагрев аккрецирующей материи в непосредственной близости от горизонта событий. Данный механизм полностью подтверждается численными расчетами магнитогидродинамики в рамках Общей теории относительности [15–21].

Современная точность Event Horizon Telescope недостаточна для проверки модифицированных теорий гравитации. В обозримом будущем такая проверка станет возможной при реализации международного проекта космической обсерватории «Миллиметрон», чье угловое разрешение будет на несколько порядков выше, чем у Event Horizon Telescope [22].

На рис. 1 показаны смоделированные темные пятна в центральной части изображения сверхмассивной черной дыры SgrA* при различных значениях параметра вращения черной дыры. Внешние замкнутые кривые соответствуют границам классических теней черных дыр, спроецированных на небесную сферу.

Соответственно, на рис. 2 представлена суперпозиция изображений сверхмассивных черных дыр SgrA* и M87*, полученных Event Horizon Telescope, со смоделированными центральными темными пятнами для значения параметра вращения a = 0,75 в обоих случаях.

Подпись к рис. 1:
Смоделированные темные пятна (черные области) в центральной части изображения сверхмассивной черной дыры SgrA* для значений параметра вращения a = 0 (левая панель), a = 0,75 (средняя панель) и a = 0,9982 (правая панель). Пунктирные кольца соответствуют воображаемым размерам горизонта черной дыры в евклидовом пространстве при отсутствии гравитации. Стрелки показывают проекцию оси вращения черной дыры на небесную сферу.

Подпись к рис. 2:
Суперпозиция изображений сверхмассивных черных дыр SgrA* (левая панель) и M87* (правая панель), полученных Event Horizon Telescope, со смоделированными формами темных пятен (черная область) для параметра вращения a = 0,75 в обоих случаях. Замкнутая кривая со «звездой» соответствует наиболее яркой части тонкого аккреционного диска, наблюдаемого удаленным наблюдателем.

Излучение внутри фотонных сфер

\begin{figure}[ht]
    \centerfloat{
        \hfill
        \subcaptionbox[List-of-Figures entry]{3D изображение SgrA* и ее фотонные траектории\label{fig:Shad1_a}}{
            \includegraphics[width=0.33\linewidth]{Shad1_a}}
        \hfill
        \subcaptionbox{Фотонная траектория с прицельными параметрами \(\lambda =  -1.493, q = 3.629\) \label{fig:Shad1_b}}
        {
        \includegraphics[width=0.33\linewidth]{Shad1_b}
        }
        \hfill
        \subcaptionbox{Граница темного пятна под углом \(\theta_0 = 82.2^\circ\)\label{fig:Shad1_2d}}
        {
            \includegraphics[width=0.33\linewidth]{Shad1_2d}
        }
        \hfill
    }
    \caption[Фотонные траектории и изображение темного пятна SgrA*]{Фотонные траектории и изображение темного пятна SgrA*}\label{fig:Shad0}
\end{figure}

\begin{figure}[ht]
    \centerfloat{
        \hfill
        \subcaptionbox[List-of-Figures entry]{3D изображение M87* и ее фотонные траектории\label{fig:Shad2_a}}{
            \includegraphics[width=0.3\linewidth]{Shad2_a}}
        \hfill
        \subcaptionbox{Граница темного пятна при \(a = 1\)\label{fig:Shad2_2d_a}}{
            \includegraphics[width=0.3\linewidth]{Shad2_2d_a}}
        \hfill
        \subcaptionbox{Граница темного пятна при \(a = 0.75\) \label{fig:Shad2_2d_b}}{
            \includegraphics[width=0.3\linewidth]{Shad2_2d_b}}
        \hfill
    }
    \caption[Фотонные траектории и изображение темного пятна M87]{Фотонные траектории и изображение темного пятна M87}\label{fig:Shad0}
\end{figure}

\begin{figure}[ht]
    \centerfloat{
        \hfill
        \subcaptionbox[List-of-Figures entry]{SgrA*\label{fig:Shad1_r}}{
            \includegraphics[width=0.3\linewidth]{Shad1_r}}
        \hfill
        \subcaptionbox{M87\label{fig:Shad2_r}}{
            \includegraphics[width=0.3\linewidth]{Shad2_r}}
        \hfill
    }
    \caption[Cуперпозиция изображений сверхмассивных черных дыр SgrA* и M87]{Cуперпозиция изображений сверхмассивных черных дыр SgrA* и M87, полученных Event Horizon Telescope, со смоделированными центральными темными пятнами для значения параметра вращения a = 0,75 в обоих случаях}\label{fig:Shad0}
\end{figure}

%Сошлёмся на библиографию.
%Одна ссылка: \cite[с.~54]{Sokolov}\cite[с.~36]{Gaidaenko}.
%Две ссылки: \cite{Sokolov,Gaidaenko}.
%Ссылка на собственные работы: \cite{vakbib1, confbib2}.
%Много ссылок: %\cite[с.~54]{Lermontov,Management,Borozda} % такой «фокус»
%%вызывает biblatex warning относительно опции sortcites, потому что неясно, к
%%какому источнику относится уточнение о страницах, а bibtex об этой проблеме
%%даже не предупреждает
%\cite{Lermontov, Management, Borozda, Marketing, Constitution, FamilyCode,
%    Gost.7.0.53, Razumovski, Lagkueva, Pokrovski, Methodology, Berestova,
%    Kriger}%
%\ifnumequal{\value{bibliosel}}{0}{% Примеры для bibtex8
%    \cite{Sirotko, Lukina, Encyclopedia, Nasirova}%
%}{% Примеры для biblatex через движок biber
%    \cite{Sirotko2, Lukina2, Encyclopedia2, Nasirova2}%
%}%
%.
%И~ещё немного ссылок:~\cite{Article,Book,Booklet,Conference,Inbook,Incollection,Manual,Mastersthesis,
%    Misc,Phdthesis,Proceedings,Techreport,Unpublished}
%% Следует обратить внимание, что пробел после запятой внутри \cite{}
%% обрабатывается ожидаемо, а пробел перед запятой, может вызывать проблемы при
%% обработке ссылок.
%\cite{medvedev2006jelektronnye, CEAT:CEAT581, doi:10.1080/01932691.2010.513279,
%    Gosele1999161,Li2007StressAnalysis, Shoji199895, test:eisner-sample,
%    test:eisner-sample-shorted, AB_patent_Pomerantz_1968, iofis_patent1960}%
%\ifnumequal{\value{bibliosel}}{0}{% Примеры для bibtex8
%}{% Примеры для biblatex через движок biber
%    \cite{patent2h, patent3h, patent2}%
%}%
%.
%
%\ifnumequal{\value{bibliosel}}{0}{% Примеры для bibtex8
%Попытка реализовать несколько ссылок на конкретные страницы
%для \texttt{bibtex} реализации библиографии:
%[\citenum{Sokolov}, с.~54; \citenum{Gaidaenko}, с.~36].
%}{% Примеры для biblatex через движок biber
%Несколько источников (мультицитата):
%% Тут специально написано по-разному тире, для демонстрации, что
%% применение специальных тире в настоящий момент в biblatex приводит к непоказу
%% "с.".
%\cites[vii--x, 5, 7]{Sokolov}[v"--~x, 25, 526]{Gaidaenko}[vii--x, 5, 7]{Techreport},
%работает только в \texttt{biblatex} реализации библиографии.
%}%
%
%Ссылки на собственные работы:~\cite{vakbib1, confbib1}.
%
%Сошлёмся на приложения: Приложение~\cref{app:A}, Приложение~\cref{app:B2}.
%
%Сошлёмся на формулу: формула~\cref{eq:equation1}.
%
%Сошлёмся на изображение: рисунок~\cref{fig:knuth}.
%
%Стандартной практикой является добавление к ссылкам префикса, характеризующего тип элемента.
%Это не является строгим требованием, но~позволяет лучше ориентироваться в документах большого размера.
%Например, для ссылок на~рисунки используется префикс \textit{fig},
%для ссылки на~таблицу "--- \textit{tab}.
%
%В таблице \cref{tab:tab_pref} приложения~\cref{app:B4} приведён список рекомендуемых
%к использованию стандартных префиксов.
%
%В некоторых ситуациях возникает необходимость отойти от требований ГОСТ по оформлению ссылок на
%литературу.
%В таком случае можно воспользоваться дополнительными опциями пакета \verb+biblatex+.
%
%Например, в ссылке на книгу~\cite{sobenin_kdv} использование опции \verb+maxnames=4+ позволяет
%вывести имена всех четырёх авторов.
%По ГОСТ имена последних трёх авторов опускаются.
%
%Кроме того, часто возникают проблемы с транслитерованными инициалами. Некоторые буквы русского
%алфавита по правилам транслитерации записываются двумя буквами латинского алфавита (ю-yu, ё-yo и
%т.д.).
%Такие инициалы \verb+biblatex+ будет сокращать до одной буквы, что неверно.
%Поправить его работу можно использовав опцию \verb+giveninits=false+.
%Пример использования этой опции можно видеть в ссылке~\cite{initials}.



%\section{Формулы}\label{sec:ch1/sec3}
%
%Благодаря пакету \textit{icomma}, \LaTeX~одинаково хорошо воспринимает
%в~качестве десятичного разделителя и запятую (\(3,1415\)), и точку (\(3.1415\)).
%
%\subsection{Ненумерованные одиночные формулы}\label{subsec:ch1/sec3/sub1}
%
%Вот так может выглядеть формула, которую необходимо вставить в~строку
%по~тексту: \(x \approx \sin x\) при \(x \to 0\).
%
%А вот так выглядит ненумерованная отдельностоящая формула c подстрочными
%и надстрочными индексами:
%\[
%    (x_1+x_2)^2 = x_1^2 + 2 x_1 x_2 + x_2^2
%\]
%
%Формула с неопределенным интегралом:
%\[
%    \int f(\alpha+x)=\sum\beta
%\]
%
%При использовании дробей формулы могут получаться очень высокие:
%\[
%    \frac{1}{\sqrt{2}+
%        \displaystyle\frac{1}{\sqrt{2}+
%            \displaystyle\frac{1}{\sqrt{2}+\cdots}}}
%\]
%
%В формулах можно использовать греческие буквы:
%%Все \original... команды заранее, ради этого примера, определены в Dissertation\userstyles.tex
%\[
%    \alpha\beta\gamma\delta\originalepsilon\epsilon\zeta\eta\theta%
%    \vartheta\iota\kappa\varkappa\lambda\mu\nu\xi\pi\varpi\rho\varrho%
%    \sigma\varsigma\tau\upsilon\originalphi\phi\chi\psi\omega\Gamma\Delta%
%    \Theta\Lambda\Xi\Pi\Sigma\Upsilon\Phi\Psi\Omega
%\]
%\[%https://texfaq.org/FAQ-boldgreek
%    \boldsymbol{\alpha\beta\gamma\delta\originalepsilon\epsilon\zeta\eta%
%        \theta\vartheta\iota\kappa\varkappa\lambda\mu\nu\xi\pi\varpi\rho%
%        \varrho\sigma\varsigma\tau\upsilon\originalphi\phi\chi\psi\omega\Gamma%
%        \Delta\Theta\Lambda\Xi\Pi\Sigma\Upsilon\Phi\Psi\Omega}
%\]
%
%Для добавления формул можно использовать пары \verb+$+\dots\verb+$+ и \verb+$$+\dots\verb+$$+,
%но~они считаются устаревшими.
%Лучше использовать их функциональные аналоги \verb+\(+\dots\verb+\)+ и \verb+\[+\dots\verb+\]+.
%
%\subsection{Ненумерованные многострочные формулы}\label{subsec:ch1/sec3/sub2}
%
%Вот так можно написать две формулы, не нумеруя их, чтобы знаки <<равно>> были
%строго друг под другом:
%\begin{align}
%    f_W & =  \min \left( 1, \max \left( 0, \frac{W_{soil} / W_{max}}{W_{crit}} \right)  \right), \nonumber \\
%    f_T & =  \min \left( 1, \max \left( 0, \frac{T_s / T_{melt}}{T_{crit}} \right)  \right), \nonumber
%\end{align}
%
%Выровнять систему ещё и по переменной \( x \) можно, используя окружение
%\verb|alignedat| из пакета \verb|amsmath|. Вот так:
%\[
%|x| = \left\{
%\begin{alignedat}{2}
%     &   & x, \quad & \text{eсли } x\geqslant 0 \\
%     & - & x, \quad & \text{eсли } x<0
%\end{alignedat}
%\right.
%\]
%Здесь первый амперсанд (в исходном \LaTeX\ описании формулы) означает
%выравнивание по~левому краю, второй "--- по~\( x \), а~третий "--- по~слову
%<<если>>. Команда \verb|\quad| делает большой горизонтальный пробел.
%
%Ещё вариант:
%\[
%    |x|=
%    \begin{cases}
%        \phantom{-}x, \text{если } x \geqslant 0 \\
%        -x, \text{если } x<0
%    \end{cases}
%\]
%
%Кроме того, для  нумерованных формул \verb|alignedat| делает вертикальное
%выравнивание номера формулы по центру формулы. Например, выравнивание
%компонент вектора:
%\begin{equation}
%    \label{eq:2p3}
%    \begin{alignedat}{2}
%        {\mathbf{N}}_{o1n}^{(j)} = \,{\sin} \phi\,n\!\left(n+1\right)
%        {\sin}\theta\,
%        \pi_n\!\left({\cos} \theta\right)
%        \frac{
%        z_n^{(j)}\!\left( \rho \right)
%        }{\rho}\,
%         & {\boldsymbol{\hat{\mathrm e}}}_{r}\,+      \\
%        +\,
%        {\sin} \phi\,
%        \tau_n\!\left({\cos} \theta\right)
%        \frac{
%        \left[\rho z_n^{(j)}\!\left( \rho \right)\right]^{\prime}
%        }{\rho}\,
%         & {\boldsymbol{\hat{\mathrm e}}}_{\theta}\,+ \\
%        +\,
%        {\cos} \phi\,
%        \pi_n\!\left({\cos} \theta\right)
%        \frac{
%        \left[\rho z_n^{(j)}\!\left( \rho \right)\right]^{\prime}
%        }{\rho}\,
%         & {\boldsymbol{\hat{\mathrm e}}}_{\phi}\:.
%    \end{alignedat}
%\end{equation}
%
%Ещё об отступах. Иногда для лучшей <<читаемости>> формул полезно
%немного исправить стандартные интервалы \LaTeX\ с учётом логической
%структуры самой формулы. Например в формуле~\cref{eq:2p3} добавлен
%небольшой отступ \verb+\,+ между основными сомножителями, ниже
%результат применения всех вариантов отступа:
%\begin{align*}
%    \backslash!             & \quad f(x) = x^2\! +3x\! +2         \\
%    \mbox{по-умолчанию}     & \quad f(x) = x^2+3x+2               \\
%    \backslash,             & \quad f(x) = x^2\, +3x\, +2         \\
%    \backslash{:}           & \quad f(x) = x^2\: +3x\: +2         \\
%    \backslash;             & \quad f(x) = x^2\; +3x\; +2         \\
%    \backslash \mbox{space} & \quad f(x) = x^2\ +3x\ +2           \\
%    \backslash \mbox{quad}  & \quad f(x) = x^2\quad +3x\quad +2   \\
%    \backslash \mbox{qquad} & \quad f(x) = x^2\qquad +3x\qquad +2
%\end{align*}
%
%Можно использовать разные математические алфавиты:
%\begin{align}
%    \mathcal{ABCDEFGHIJKLMNOPQRSTUVWXYZ} \nonumber  \\
%    \mathfrak{ABCDEFGHIJKLMNOPQRSTUVWXYZ} \nonumber \\
%    \mathbb{ABCDEFGHIJKLMNOPQRSTUVWXYZ} \nonumber
%\end{align}
%
%Посмотрим на систему уравнений на примере аттрактора Лоренца:
%
%\[
%\left\{
%\begin{array}{rl}
%    \dot x = & \sigma (y-x)  \\
%    \dot y = & x (r - z) - y \\
%    \dot z = & xy - bz
%\end{array}
%\right.
%\]
%
%А для вёрстки матриц удобно использовать многоточия:
%\[
%    \left(
%        \begin{array}{ccc}
%            a_{11} & \ldots & a_{1n} \\
%            \vdots & \ddots & \vdots \\
%            a_{n1} & \ldots & a_{nn} \\
%        \end{array}
%    \right)
%\]
%
%\subsection{Нумерованные формулы}\label{subsec:ch1/sec3/sub3}
%
%А вот так пишется нумерованная формула:
%\begin{equation}
%    \label{eq:equation1}
%    e = \lim_{n \to \infty} \left( 1+\frac{1}{n} \right) ^n
%\end{equation}
%
%Нумерованных формул может быть несколько:
%\begin{equation}
%    \label{eq:equation2}
%    \lim_{n \to \infty} \sum_{k=1}^n \frac{1}{k^2} = \frac{\pi^2}{6}
%\end{equation}
%
%Впоследствии на формулы~\cref{eq:equation1, eq:equation2} можно ссылаться.
%
%Сделать так, чтобы номер формулы стоял напротив средней строки, можно,
%используя окружение \verb|multlined| (пакет \verb|mathtools|) вместо
%\verb|multline| внутри окружения \verb|equation|. Вот так:
%\begin{equation} % \tag{S} % tag - вписывает свой текст
%    \label{eq:equation3}
%    \begin{multlined}
%        1+ 2+3+4+5+6+7+\dots + \\
%        + 50+51+52+53+54+55+56+57 + \dots + \\
%        + 96+97+98+99+100=5050
%    \end{multlined}
%\end{equation}
%
%Уравнения~\cref{eq:subeq_1,eq:subeq_2} демонстрируют возможности
%окружения \verb|subequations| (пакет \verb|amsmath|).
%\begin{subequations}
%    \label{eq:subeq_1}
%    \begin{gather}
%        y = x^2 + 1 \label{eq:subeq_1-1} \\
%        y = 2 x^2 - x + 1 \label{eq:subeq_1-2}
%    \end{gather}
%\end{subequations}
%Ссылки на отдельные уравнения~\cref{eq:subeq_1-1,eq:subeq_1-2,eq:subeq_2-1}.
%\begin{subequations}
%    \label{eq:subeq_2}
%    \begin{align}
%        y & = x^3 + x^2 + x + 1 \label{eq:subeq_2-1} \\
%        y & = x^2
%    \end{align}
%\end{subequations}
%
%\subsection{Форматирование чисел и размерностей величин}\label{sec:units}
%
%Числа форматируются при помощи команды \verb|\num|:
%\num{5,3};
%\num{2,3e8};
%\num{12345,67890};
%\num{2,6 d4};
%\num{1+-2i};
%\num{.3e45};
%\num[exponent-base=2]{5 e64};
%\num[exponent-base=2,exponent-to-prefix]{5 e64};
%\num{1.654 x 2.34 x 3.430}
%\num{1 2 x 3 / 4}.
%Для написания последовательности чисел можно использовать команды \verb|\numlist| и \verb|\numrange|:
%\numlist{10;30;50;70}; \numrange{10}{30}.
%Значения углов можно форматировать при помощи команды \verb|\ang|:
%\ang{2.67};
%\ang{30,3};
%\ang{-1;;};
%\ang{;-2;};
%\ang{;;-3};
%\ang{300;10;1}.
%
%Обратите внимание, что ГОСТ запрещает использование знака <<->> для обозначения отрицательных чисел
%за исключением формул, таблиц и~рисунков.
%Вместо него следует использовать слово <<минус>>.
%
%Размерности можно записывать при помощи команд \verb|\si| и \verb|\SI|:
%\si{\farad\squared\lumen\candela};
%\si{\joule\per\mole\per\kelvin};
%\si[per-mode = symbol-or-fraction]{\joule\per\mole\per\kelvin};
%\si{\metre\per\second\squared};
%\SI{0.10(5)}{\neper};
%\SI{1.2-3i e5}{\joule\per\mole\per\kelvin};
%\SIlist{1;2;3;4}{\tesla};
%\SIrange{50}{100}{\volt}.
%Список единиц измерений приведён в таблицах~\cref{tab:unit:base,
%    tab:unit:derived,tab:unit:accepted,tab:unit:physical,tab:unit:other}.
%Приставки единиц приведены в~таблице~\cref{tab:unit:prefix}.
%
%С дополнительными опциями форматирования можно ознакомиться в~описании пакета \texttt{siunitx};
%изменить или добавить единицы измерений можно в~файле \texttt{siunitx.cfg}.
%
%\begin{table}
%    \centering
%    \captionsetup{justification=centering} % выравнивание подписи по-центру
%    \caption{Основные величины СИ}\label{tab:unit:base}
%    \begin{tabular}{llc}
%        \toprule
%        Название  & Команда          & Символ         \\
%        \midrule
%        Ампер     & \verb|\ampere|   & \si{\ampere}   \\
%        Кандела   & \verb|\candela|  & \si{\candela}  \\
%        Кельвин   & \verb|\kelvin|   & \si{\kelvin}   \\
%        Килограмм & \verb|\kilogram| & \si{\kilogram} \\
%        Метр      & \verb|\metre|    & \si{\metre}    \\
%        Моль      & \verb|\mole|     & \si{\mole}     \\
%        Секунда   & \verb|\second|   & \si{\second}   \\
%        \bottomrule
%    \end{tabular}
%\end{table}
%
%\begin{table}
%    \small
%    \centering
%    \begin{threeparttable}% выравнивание подписи по границам таблицы
%        \caption{Производные единицы СИ}\label{tab:unit:derived}
%        \begin{tabular}{llc|llc}
%            \toprule
%            Название       & Команда               & Символ              & Название & Команда & Символ \\
%            \midrule
%            Беккерель      & \verb|\becquerel|     & \si{\becquerel}     &
%            Ньютон         & \verb|\newton|        & \si{\newton}                                      \\
%            Градус Цельсия & \verb|\degreeCelsius| & \si{\degreeCelsius} &
%            Ом             & \verb|\ohm|           & \si{\ohm}                                         \\
%            Кулон          & \verb|\coulomb|       & \si{\coulomb}       &
%            Паскаль        & \verb|\pascal|        & \si{\pascal}                                      \\
%            Фарад          & \verb|\farad|         & \si{\farad}         &
%            Радиан         & \verb|\radian|        & \si{\radian}                                      \\
%            Грей           & \verb|\gray|          & \si{\gray}          &
%            Сименс         & \verb|\siemens|       & \si{\siemens}                                     \\
%            Герц           & \verb|\hertz|         & \si{\hertz}         &
%            Зиверт         & \verb|\sievert|       & \si{\sievert}                                     \\
%            Генри          & \verb|\henry|         & \si{\henry}         &
%            Стерадиан      & \verb|\steradian|     & \si{\steradian}                                   \\
%            Джоуль         & \verb|\joule|         & \si{\joule}         &
%            Тесла          & \verb|\tesla|         & \si{\tesla}                                       \\
%            Катал          & \verb|\katal|         & \si{\katal}         &
%            Вольт          & \verb|\volt|          & \si{\volt}                                        \\
%            Люмен          & \verb|\lumen|         & \si{\lumen}         &
%            Ватт           & \verb|\watt|          & \si{\watt}                                        \\
%            Люкс           & \verb|\lux|           & \si{\lux}           &
%            Вебер          & \verb|\weber|         & \si{\weber}                                       \\
%            \bottomrule
%        \end{tabular}
%    \end{threeparttable}
%\end{table}
%
%\begin{table}
%    \centering
%    \begin{threeparttable}% выравнивание подписи по границам таблицы
%        \caption{Внесистемные единицы}\label{tab:unit:accepted}
%
%        \begin{tabular}{llc}
%            \toprule
%            Название        & Команда           & Символ          \\
%            \midrule
%            День            & \verb|\day|       & \si{\day}       \\
%            Градус          & \verb|\degree|    & \si{\degree}    \\
%            Гектар          & \verb|\hectare|   & \si{\hectare}   \\
%            Час             & \verb|\hour|      & \si{\hour}      \\
%            Литр            & \verb|\litre|     & \si{\litre}     \\
%            Угловая минута  & \verb|\arcminute| & \si{\arcminute} \\
%            Угловая секунда & \verb|\arcsecond| & \si{\arcsecond} \\ %
%            Минута          & \verb|\minute|    & \si{\minute}    \\
%            Тонна           & \verb|\tonne|     & \si{\tonne}     \\
%            \bottomrule
%        \end{tabular}
%    \end{threeparttable}
%\end{table}
%
%\begin{table}
%    \centering
%    \captionsetup{justification=centering}
%    \caption{Внесистемные единицы, получаемые из эксперимента}\label{tab:unit:physical}
%    \begin{tabular}{llc}
%        \toprule
%        Название                & Команда                  & Символ                 \\
%        \midrule
%        Астрономическая единица & \verb|\astronomicalunit| & \si{\astronomicalunit} \\
%        Атомная единица массы   & \verb|\atomicmassunit|   & \si{\atomicmassunit}   \\
%        Боровский радиус        & \verb|\bohr|             & \si{\bohr}             \\
%        Скорость света          & \verb|\clight|           & \si{\clight}           \\
%        Дальтон                 & \verb|\dalton|           & \si{\dalton}           \\
%        Масса электрона         & \verb|\electronmass|     & \si{\electronmass}     \\
%        Электрон Вольт          & \verb|\electronvolt|     & \si{\electronvolt}     \\
%        Элементарный заряд      & \verb|\elementarycharge| & \si{\elementarycharge} \\
%        Энергия Хартри          & \verb|\hartree|          & \si{\hartree}          \\
%        Постоянная Планка       & \verb|\planckbar|        & \si{\planckbar}        \\
%        \bottomrule
%    \end{tabular}
%\end{table}
%
%\begin{table}
%    \centering
%    \begin{threeparttable}% выравнивание подписи по границам таблицы
%        \caption{Другие внесистемные единицы}\label{tab:unit:other}
%        \begin{tabular}{llc}
%            \toprule
%            Название                  & Команда              & Символ             \\
%            \midrule
%            Ангстрем                  & \verb|\angstrom|     & \si{\angstrom}     \\
%            Бар                       & \verb|\bar|          & \si{\bar}          \\
%            Барн                      & \verb|\barn|         & \si{\barn}         \\
%            Бел                       & \verb|\bel|          & \si{\bel}          \\
%            Децибел                   & \verb|\decibel|      & \si{\decibel}      \\
%            Узел                      & \verb|\knot|         & \si{\knot}         \\
%            Миллиметр ртутного столба & \verb|\mmHg|         & \si{\mmHg}         \\
%            Морская миля              & \verb|\nauticalmile| & \si{\nauticalmile} \\
%            Непер                     & \verb|\neper|        & \si{\neper}        \\
%            \bottomrule
%        \end{tabular}
%    \end{threeparttable}
%\end{table}
%
%\begin{table}
%    \small
%    \centering
%    \begin{threeparttable}% выравнивание подписи по границам таблицы
%        \caption{Приставки СИ}\label{tab:unit:prefix}
%        \begin{tabular}{llcc|llcc}
%            \toprule
%            Приставка & Команда       & Символ      & Степень      &
%            Приставка & Команда       & Символ      & Степень        \\
%            \midrule
%            Иокто     & \verb|\yocto| & \si{\yocto} & \textminus24 &
%            Дека      & \verb|\deca|  & \si{\deca}  & 1              \\
%            Зепто     & \verb|\zepto| & \si{\zepto} & \textminus21 &
%            Гекто     & \verb|\hecto| & \si{\hecto} & 2              \\
%            Атто      & \verb|\atto|  & \si{\atto}  & \textminus18 &
%            Кило      & \verb|\kilo|  & \si{\kilo}  & 3              \\
%            Фемто     & \verb|\femto| & \si{\femto} & \textminus15 &
%            Мега      & \verb|\mega|  & \si{\mega}  & 6              \\
%            Пико      & \verb|\pico|  & \si{\pico}  & \textminus12 &
%            Гига      & \verb|\giga|  & \si{\giga}  & 9              \\
%            Нано      & \verb|\nano|  & \si{\nano}  & \textminus9  &
%            Терра     & \verb|\tera|  & \si{\tera}  & 12             \\
%            Микро     & \verb|\micro| & \si{\micro} & \textminus6  &
%            Пета      & \verb|\peta|  & \si{\peta}  & 15             \\
%            Милли     & \verb|\milli| & \si{\milli} & \textminus3  &
%            Екса      & \verb|\exa|   & \si{\exa}   & 18             \\
%            Санти     & \verb|\centi| & \si{\centi} & \textminus2  &
%            Зетта     & \verb|\zetta| & \si{\zetta} & 21             \\
%            Деци      & \verb|\deci|  & \si{\deci}  & \textminus1  &
%            Иотта     & \verb|\yotta| & \si{\yotta} & 24             \\
%            \bottomrule
%        \end{tabular}
%    \end{threeparttable}
%\end{table}
%
%\subsection{Заголовки с формулами: \texorpdfstring{\(a^2 + b^2 = c^2\)}{%
%        a\texttwosuperior\ + b\texttwosuperior\ = c\texttwosuperior},
%    \texorpdfstring{\(\left\vert\textrm{{Im}}\Sigma\left(
%            \protect\varepsilon\right)\right\vert\approx const\)}{|ImΣ (ε)| ≈ const},
%    \texorpdfstring{\(\sigma_{xx}^{(1)}\)}{σ\_\{xx\}\textasciicircum\{(1)\}}
%}\label{subsec:with_math}
%
%Пакет \texttt{hyperref} берёт текст для закладок в pdf-файле из~аргументов
%команд типа \verb|\section|, которые могут содержать математические формулы,
%а~также изменения цвета текста или шрифта, которые не отображаются в~закладках.
%Чтобы использование формул в заголовках не вызывало в~логе компиляции появление
%предупреждений типа <<\texttt{Token not allowed in~a~PDF string
%    (Unicode):(hyperref) removing...}>>, следует использовать конструкцию
%\verb|\texorpdfstring{}{}|, где в~первых фигурных скобках указывается
%формула, а~во~вторых "--- запись формулы для закладок.
%
%\section{Рецензирование текста}\label{sec:markup}
%
%В шаблоне для диссертации и автореферата заданы команды рецензирования.
%Они видны при компиляции шаблона в режиме черновика или при установке
%соответствующей настройки (\verb+showmarkup+) в~файле \verb+common/setup.tex+.
%
%Команда \verb+\todo+ отмечает текст красным цветом.
%\todo{Например, так.}
%
%Команда \verb+\note+ позволяет выбрать цвет текста.
%\note{Чёрный, } \note[red]{красный, } \note[green]{зелёный, }
%\note[blue]{синий.} \note[orange]{Обратите внимание на ширину и расстановку
%    формирующихся пробелов, в~результате приведённой записи (зависит также
%    от~применяемого компилятора).}
%
%Окружение \verb+commentbox+ также позволяет выбрать цвет.
%
%\begin{commentbox}[red]
%    Красный текст.
%
%    Несколько параграфов красного текста.
%\end{commentbox}
%
%\begin{commentbox}[blue]
%    Синяя формула.
%
%    \begin{equation}
%        \alpha + \beta = \gamma
%    \end{equation}
%\end{commentbox}
%
%\verb+commentbox+ позволяет закомментировать участок кода в~режиме чистовика.
%Чтобы убрать кусок кода для всех режимов, можно использовать окружение
%\verb+comment+.
%
%\begin{comment}
%Этот текст всегда скрыт.
%\end{comment}
%
%\section{Работа со списком сокращений и~условных обозначений}\label{sec:acronyms}
%
%С помощью пакета \texttt{nomencl} можно создавать удобный сортированный список
%сокращений и условных обозначений во время написания текста. Вызов
%\verb+\nomenclature+ добавляет нужный символ или сокращение с~описанием
%в~список, который затем печатается вызовом \verb+\printnomenclature+
%в~соответствующем разделе.

%makeindex -s nomencl.ist -o dissertation.nls dissertation.nlo

%Для того, чтобы эти операции прошли, потребуется дополнительный вызов
%\verb+makeindex -s nomencl.ist -o %.nls %.nlo+ в~командной строке, где вместо
%\verb+%+ следует подставить имя главного файла проекта (\verb+dissertation+
%для этого шаблона).
%Затем потребуется один или два дополнительных вызова компилятора проекта.
%\begin{equation}
%    \omega = c k,
%\end{equation}
%где \( \omega \) "--- частота света, \( c \) "--- скорость света, \( k \) "---
%модуль волнового вектора.
%\nomenclature{\(\omega\)}{частота света\nomrefeq}
%\nomenclature{\(c\)}{скорость света\nomrefpage}
%\nomenclature{\(k\)}{модуль волнового вектора\nomrefeqpage}
%Использование
%\begin{verbatim}
%\nomenclature{\(\omega\)}{частота света\nomrefeq}
%\nomenclature{\(c\)}{скорость света\nomrefpage}
%\nomenclature{\(k\)}{модуль волнового вектора\nomrefeqpage}
%\end{verbatim}
%после уравнения добавит в список условных обозначений три записи.
%Ссылки \verb+\nomrefeq+ на последнее уравнение, \verb+\nomrefpage+ "--- на
%страницу, \verb+\nomrefeqpage+ "--- сразу на~последнее уравнение и~на~страницу,
%можно опускать и~не~использовать.
%
%Группировкой и сортировкой пунктов в списке можно управлять с~помощью указания
%дополнительных аргументов к команде \verb+nomenclature+.
%Например, при вызове
%\begin{verbatim}
%\nomenclature[03]{\( \hbar \)}{постоянная Планка}
%\nomenclature[01]{\( G \)}{гравитационная постоянная}
%\end{verbatim}
%\( G \) будет стоять в списке выше, чем \( \hbar \).
%Для корректных вертикальных отступов между строками в описании лучше
%не~использовать многострочные формулы в~списке обозначений.
%
%\nomenclature{%
%    \( \begin{rcases}
%        a_n \\
%        b_n
%    \end{rcases} \)%
%}{коэффициенты разложения Ми в дальнем поле соответствующие электрическим и
%    магнитным мультиполям}
%\nomenclature[a\( e \)]{\( {\boldsymbol{\hat{\mathrm e}}} \)}{единичный вектор}
%\nomenclature{\( E_0 \)}{амплитуда падающего поля}
%\nomenclature{\( j \)}{тип функции Бесселя}
%\nomenclature{\( k \)}{волновой вектор падающей волны}
%\nomenclature{%
%    \( \begin{rcases}
%        a_n \\
%        b_n
%    \end{rcases} \)%
%}{и снова коэффициенты разложения Ми в дальнем поле соответствующие
%    электрическим и магнитным мультиполям. Добавлено много текста, так что
%    описание группы условных обозначений значительно превысило высоту этой
%    группы...}
%\nomenclature{\( L \)}{общее число слоёв}
%\nomenclature{\( l \)}{номер слоя внутри стратифицированной сферы}
%\nomenclature{\( \lambda \)}{длина волны электромагнитного излучения в вакууме}
%\nomenclature{\( n \)}{порядок мультиполя}
%\nomenclature{%
%    \( \begin{rcases}
%        {\mathbf{N}}_{e1n}^{(j)} & {\mathbf{N}}_{o1n}^{(j)} \\
%        {\mathbf{M}_{o1n}^{(j)}} & {\mathbf{M}_{e1n}^{(j)}}
%    \end{rcases} \)%
%}{сферические векторные гармоники}
%\nomenclature{\( \mu \)}{магнитная проницаемость в вакууме}
%\nomenclature{\( r, \theta, \phi \)}{полярные координаты}
%\nomenclature{\( \omega \)}{частота падающей волны}
%
%С помощью \verb+nomenclature+ можно включать в~список сокращения,
%не~используя их~в~тексте.
%% запись сокращения в список происходит командой \nomenclature,
%% а не употреблением самого сокращения
%\nomenclature{FEM}{finite element method, метод конечных элементов}
%\nomenclature{FIT}{finite integration technique, метод конечных интегралов}
%\nomenclature{FMM}{fast multipole method, быстрый метод многополюсника}
%\nomenclature{FVTD}{finite volume time-domain, метод конечных объёмов
%    во~временной области}
%\nomenclature{MLFMA}{multilevel fast multipole algorithm, многоуровневый
%    быстрый алгоритм многополюсника}
%\nomenclature{BEM}{boundary element method, метод граничных элементов}
%\nomenclature{CST MWS}{Computer Simulation Technology Microwave Studio
%    программа для компьютерного моделирования уравнен Максвелла}
%\nomenclature{DDA}{discrete dipole approximation, приближение дискретных
%    диполей}
%\nomenclature{FDFD}{finite difference frequency domain, метод конечных
%    разностей в~частотной области}
%\nomenclature{FDTD}{finite difference time domain, метод конечных разностей
%    во~временной области}
%\nomenclature{MoM}{method of moments, метод моментов}
%\nomenclature{MSTM}{multiple sphere T-Matrix, метод Т-матриц для множества
%    сфер}
%\nomenclature{PSTD}{pseudospectral time domain method, псевдоспектральный метод
%    во~временной области}
%\nomenclature{TLM}{transmission line matrix method, метод матриц линий передач}

\FloatBarrier
